\documentclass[aspectratio=1610]{beamer}
\usepackage[utf8]{inputenc}
\usepackage[T1]{fontenc}
%%%%%%%
\usepackage{hyperref}
\usepackage{mathtools}
% \usepackage{layout}
% \usepackage{lipsum}
%%%%%%%
\usefonttheme{professionalfonts} % keep beamer from replacing the math fonts
\usetheme[%%%%%%%%% Complete settings. Default (without argument) value in []
% titleimagecolor=red,                 % [gray], darkgray, red, blue, green
% titleimagemargin=2mm,           % Distance [2mm]    Frame around title page image
% navigationsymbols=false,         % true   / [false]  Navigation symbols in the foot
% mathseriffont=false,                 % true   / [false]  Serif / non-serif math fonts
% foot=true,                                 % true / [false]    Footline or not
% slidenumnofoot=true,              % [true]   / false  Keep slide num even when foot=false
% blackenumeratenumber=true,  % [true] / false    Black enumerate numbers, o.w. Lund bronze
% blackitemmark=false,              % true   / [false]  Black item marks, o.w. Lund bronze
% defaultfont=false,                    % true   / [false]  Falls back to default beamer fonts
% sectionframe=true,
% centeralign = true,                   % true/[false]  Vertically alignes text on the slide. Othewise top aligned (default)
% titleline = true,                         % true/[false] line to separete the title from the text
%%%%%%%%%% Logo version
%footlogo=false,                         % [true] / false    Put LU logo to the right in footer (LU style)
%logo=LTH,
%logolang = swe,                        %  swe/[eng] language for logo
%logocol=BW,                             % RGB/BW/neg  
% LUseal = false,                         % [true]/false LU sigill i nedre högra hörnet
%%%%%%%%%%% Template variants
%%%%%%%%%%%%% LTH
% LTHbg = true,                          % true/[false]  Gray backgruond on slides (in LTH template)
% LTHtemplate,                           % LTH version of template, overrides and sets logo=LTH, footlogo=false, LUseal=false, LTHbg=true
]{ulund}
%%%% Math in Computer Modern (standard LaTeX) instead of the theme's Palatino
\DeclareSymbolFont{operators}{OT1}{cmr}{m}{n}
\DeclareSymbolFont{letters}{OML}{cmm}{m}{it}
\DeclareSymbolFont{symbols}{OMS}{cmsy}{m}{n}
\DeclareSymbolFont{largesymbols}{OMX}{cmex}{m}{n}
\SetSymbolFont{operators}{bold}{OT1}{cmr}{bx}{n}
\SetSymbolFont{letters}{bold}{OML}{cmm}{b}{it}
\SetSymbolFont{symbols}{bold}{OMS}{cmsy}{b}{n}
\DeclareMathAlphabet{\mathrm}{OT1}{cmr}{m}{n}
\DeclareMathAlphabet{\mathit}{OT1}{cmr}{m}{it}
\DeclareMathAlphabet{\mathbf}{OT1}{cmr}{bx}{n}
%%%% Draw the LU logo again on top of the slide content (the theme puts it in the background)
\makeatletter
\addtobeamertemplate{footline}{}{%
  \ifulund@printlogo
    \begin{tikzpicture}[remember picture,overlay]
      \node[anchor=south east,inner sep=0pt] at ([xshift=-5mm,yshift=4mm]current page.south east)
        {\includegraphics[height=13mm]{\ulund@logopath}};
    \end{tikzpicture}%
  \fi}
\makeatother
%%%%%%%%%%%%%%%%%%%%% Layout commands 
%%%% Foot
% \ulundfootleft{\insertshortauthor}
% \ulundfootmid{\insertshorttitle}
% \ulundfootright{\insertframenumber}% {\insertframenumber:\inserttotalframenumber}
%%%% Titleimage
%\titleimage{example-image} % Replaces the LU image.
%\titleimage{Pictures/Fikapaus}
%\titleimage{Pictures/titlepictureGroup}
%%%% Redefine LTH background color
%5 Default RGB: 216,212,207
%\definecolor{LTHgrey}{RGB}{200,200,255}
%%%%%%%%%%%%%%%%%%%%%%%%%%%%%%%%%%%
\setbeamertemplate{theorems}[numbered] % Definition 1, 2, ... as in the report
\newtheorem{proposition}{Proposition}
\newtheorem{corol}{Corollary} % own counter, numbered as in the report
\newtheorem{thm}{Theorem} % own counter, numbered as in the report
\newtheorem{lem}{Lemma} % own counter, numbered as in the report
\title[The sharp $A_p$ bound for $M$]{The sharp $A_p$ bound for the Hardy--Littlewood maximal operator}
\author[Simon Gustafsson]{Simon Gustafsson}
%%%%%%%%%%%%%%%%%%%%%
%%%%%%%%%%%%%%%%%%%%%
\begin{document}
\begin{frame}[plain]% Use plain to suppress footline box
  \titlepage
\end{frame}

%%%%%%%%%%%%%%% 
\begin{frame}
  \frametitle{Weights and weighted $L^p$ spaces}
\begin{definition}[Weights]
A \emph{weight} is a function $w \in L^1_{\mathrm{loc}}(\mathbb{R}^d)$ with $w \ge 0$ that is not almost everywhere zero.
\end{definition}
\pause
\begin{definition}[Weighted $L^p$ spaces]
Let $w$ be a weight and let $1 \le p < \infty$. The weighted $L^p$ space is
\[
  L^p_w(\mathbb{R}^d) := \left\{ f : \mathbb{R}^d \to \mathbb{R} \text{ measurable} : \int_{\mathbb{R}^d} |f(y)|^p\, w(y)\,dy < \infty \right\},
\]
with norm $\|f\|_{L^p_w} := \left( \int_{\mathbb{R}^d} |f(y)|^p\, w(y)\,dy \right)^{1/p}$. Moreover, $L^\infty_w(\mathbb{R}^d)$ is the space of measurable functions $f$ with $\|f\|_{L^\infty_w} := \inf\{ t \ge 0 : w(\{|f| > t\}) = 0 \} < \infty$.
\end{definition}
\end{frame}


%%%%%%%%%%%%%%%
\begin{frame}
  \frametitle{The Hardy--Littlewood maximal operator}
\begin{definition}[Hardy--Littlewood maximal operator]
\[
  M_w f(x) := \sup_{Q \ni x} \frac{1}{w(Q)} \int_Q |f(y)|\, w(y)\,dy .
\]
$M^c_w$ takes the supremum over centred cubes $Q(x,r)$ only, and $w \equiv 1$ gives $M$ and $M^c$.
\end{definition}
\pause
\begin{block}{$M$ bounded on $L^p_w(\mathbb{R}^d)$}
\[
  \|Mf\|_{L^p_w} = \left( \int_{\mathbb{R}^d} (Mf)^p\, w\,dy \right)^{1/p} \le C \left( \int_{\mathbb{R}^d} |f|^p\, w\,dy \right)^{1/p} = C\, \|f\|_{L^p_w}
\]
\end{block}
\end{frame}

%%%%%%%%%%%%%%%
\begin{frame}
  \frametitle{A necessary condition on $w$}
\begin{proposition}[Muckenhoupt]
If $M$ is bounded on $L^p_w(\mathbb{R}^d)$ with constant $C$, then every cube $Q$ satisfies
\[
  \left( \frac{1}{|Q|} \int_Q w\,dy \right) \left( \frac{1}{|Q|} \int_Q w^{-\frac{1}{p-1}}\,dy \right)^{p-1} \le C^p .
\]
\end{proposition}
\pause
\vspace{-0.5em}
\begin{block}{Definitions 4 and 5 ($A_p$ characteristic, dual weight)}
\vspace{-1.2em}
\[
  [w]_{A_p} := \sup_{Q} \left( \frac{1}{|Q|} \int_Q w\,dy \right) \left( \frac{1}{|Q|} \int_Q w^{-\frac{1}{p-1}}\,dy \right)^{p-1}
  = \sup_Q \frac{w(Q)}{|Q|} \left( \frac{\sigma(Q)}{|Q|} \right)^{p-1},
\]
where $\sigma := w^{-1/(p-1)}$ is the \emph{dual weight} of $w$, and $w \in A_p$ if $[w]_{A_p} < \infty$.
\end{block}
\end{frame}

%%%%%%%%%%%%%%%
\begin{frame}
  \frametitle{$A_p$ is necessary and sufficient}
\begin{corol}
If $M$ is bounded on $L^p_w(\mathbb{R}^d)$ with constant $C$, then $w \in A_p$ with
\[
  [w]_{A_p} \le C^p \quad\iff\quad [w]_{A_p}^{\frac{1}{p}} \le C .
\]
\end{corol}
\pause
\begin{thm}[Muckenhoupt 1972]
Let $1 < p < \infty$ and $w$ a weight. Then $M$ is bounded on $L^p_w(\mathbb{R}^d)$ iff $w \in A_p$.
\end{thm}
\end{frame}



%%%%%%%%%%%%%%%
\begin{frame}
  \frametitle{Main result}
\begin{thm}[Buckley 1993, Lerner 2008]
Let $1 < p < \infty$.
\begin{enumerate}[(i)]
  \item There is a constant $D_{p,d}$, depending only on $p$ and $d$, such that for every $w \in A_p$ and every $f \in L^p_w(\mathbb{R}^d)$,
  \[
    \|Mf\|_{L^p_w} \le D_{p,d}\, [w]_{A_p}^{\frac{1}{p-1}}\, \|f\|_{L^p_w} .
  \]
  \pause
  \item If $\alpha < \frac{1}{p-1}$, then there is no constant $D$ such that $\|Mf\|_{L^p_w} \le D\, [w]_{A_p}^{\alpha}\, \|f\|_{L^p_w}$ for all $w \in A_p$ and all $f \in L^p_w(\mathbb{R}^d)$.
\end{enumerate}
\end{thm}
\end{frame}

%%%%%%%%%%%%%%%
\begin{frame}
  \frametitle{Preliminaries: Besicovitch covering lemma}
\small
\setcounter{lem}{3}
\begin{lem}[Besicovitch]
Let $E \subset \mathbb{R}^d$ be bounded, and for each $x \in E$ let a centred cube $Q(x, r_x)$ be given. Then there are $x_1, x_2, \ldots \in E$ such that the cubes $Q(x_j, r_{x_j})$ cover $E$ and every point of $\mathbb{R}^d$ lies in at most $N_d$ of them.
\end{lem}
\pause
\vspace{-0.6em}
\begin{block}{Compare: Vitali covering lemma}
Let $\{B_\ell\}_{\ell=1}^N$ be a finite collection of open balls in $\mathbb{R}^d$. Then there is a finite subcollection $\{B_{\ell_j}\}_{j=1}^M$ which is pairwise disjoint and satisfies
$\big|\bigcup_{\ell=1}^N B_\ell\big| \le 3^d \big|\bigcup_{j=1}^M B_{\ell_j}\big|$.
\end{block}
\pause
\vspace{-0.6em}
\begin{block}{Choice of the cubes}
Choose $x_1 \in E$ with $r_{x_1} > \frac12 \sup_{x \in E} r_x$. Given $x_1, \ldots, x_m$, choose an uncovered point $x_{m+1}$ with
\vspace{-0.4em}
\[
  r_{x_{m+1}} > \tfrac12 \sup\Big\{ r_x : x \in E \setminus \bigcup_{j \le m} Q(x_j, r_{x_j}) \Big\}.
\]
\end{block}
\end{frame}

%%%%%%%%%%%%%%%
\begin{frame}
  \frametitle{Preliminaries: weak type $(1,1)$ for $M^c_v$}
\begin{lem}
Let $v$ be a weight with $v(Q) > 0$ for every cube $Q$. Then for every $t > 0$,
\[
  v\big(\{M^c_v f > t\}\big) \le \tfrac{N_d}{t}\, \|f\|_{L^1_v} .
\]
\end{lem}
\pause
\vspace{-0.8em}
\begin{block}{Key estimate}
\small
For $E_k := \{|x|_\infty < k,\ M^c_v f(x) > t\}$, choose $r_x$ with average $> t$ and apply Lemma 4:
\vspace{-0.4em}
\begin{align*}
  v(E_k) &\le \sum_j v\big(Q(x_j, r_{x_j})\big) \le \frac{1}{t} \sum_j \int_{Q(x_j, r_{x_j})} |f|\, v\,dy \\
  &= \frac{1}{t} \int_{\mathbb{R}^d} \Big( \sum_j \chi_{Q(x_j, r_{x_j})} \Big) |f|\, v\,dy \le \frac{N_d}{t}\, \|f\|_{L^1_v} .
\end{align*}
\end{block}
\end{frame}

%%%%%%%%%%%%%%%
\begin{frame}
  \frametitle{Preliminaries: strong type for $M^c_v$}
\begin{lem}
Let $v$ be a weight with $v(Q) > 0$ for every cube $Q$, and let $1 < q < \infty$. Then there is a constant $B_{q,d}$, depending only on $q$ and $d$, such that
\[
  \|M^c_v f\|_{L^q_v} \le B_{q,d}\, \|f\|_{L^q_v}
\]
for all $f \in L^q_v(\mathbb{R}^d)$.
\end{lem}
\pause
\begin{block}{Outline}
As for $M$: the weak type $(1,1)$ bound of Lemma 5 and the trivial bound $\|M^c_v f\|_{L^\infty_v} \le \|f\|_{L^\infty_v}$ give the $L^q_v$ bound by Marcinkiewicz interpolation. Neither bound depends on $v$, so neither does $B_{q,d}$.
\end{block}
\end{frame}

%%%%%%%%%%%%%%%
\begin{frame}
  \frametitle{Lerner's pointwise bound}
\begin{lem}[Lerner]
Let $1 < p < \infty$, $w \in A_p$ and $f \in L^1_{\mathrm{loc}}(\mathbb{R}^d)$. Then, for every $x \in \mathbb{R}^d$,
\[
  M^c f(x) \le 3^{\frac{dp}{p-1}}\, [w]_{A_p}^{\frac{1}{p-1}} \Big( M^c_w\Big( \big(M^c_\sigma(|f|\sigma^{-1})\big)^{p-1} w^{-1} \Big)(x) \Big)^{\frac{1}{p-1}} .
\]
\end{lem}
\end{frame}

%%%%%%%%%%%%%%%
\begin{frame}
  \frametitle{Lerner's pointwise bound, Proof}
\small
\begin{block}{The $A_p$ factor}
For $x \in \mathbb{R}^d$ and $r > 0$, $Q(x,r) \subseteq Q(x,3r)$ and $|Q(x,3r)| = 3^d\,|Q(x,r)|$ give
\[
  \frac{w(Q(x,r))\,\sigma(Q(x,3r))^{p-1}}{|Q(x,r)|^{p}} \le 3^{dp}\, \frac{w(Q(x,3r))}{|Q(x,3r)|} \left( \frac{\sigma(Q(x,3r))}{|Q(x,3r)|} \right)^{p-1} \le 3^{dp}\, [w]_{A_p}.
\]
\end{block}
\pause
\begin{block}{The inner average}
Every $z \in Q(x,r)$ satisfies $Q(x,r) \subseteq Q(z,2r) \subseteq Q(x,3r)$, so
\[
  \frac{1}{\sigma(Q(x,3r))} \int_{Q(x,r)} |f|\,\sigma^{-1}\sigma\,dy \le \frac{1}{\sigma(Q(z,2r))} \int_{Q(z,2r)} |f|\,\sigma^{-1}\sigma\,dy \le M^c_\sigma\big(|f|\sigma^{-1}\big)(z).
\]
\end{block}
\end{frame}

%%%%%%%%%%%%%%%
\begin{frame}
  \frametitle{Proof of Theorem 2 (i)}
\small
\vspace{-0.8em}
\begin{align*}
  \|Mf\|_{L^p_w}
  &\le 2^d\, \| M^c f \|_{L^p_w} \\
  \onslide<2->{&\le 2^d\, \Big\| 3^{\frac{dp}{p-1}}\, [w]_{A_p}^{\frac{1}{p-1}} \Big( M^c_w\Big( \big(M^c_\sigma(|f|\sigma^{-1})\big)^{p-1} w^{-1} \Big) \Big)^{\frac{1}{p-1}} \Big\|_{L^p_w} \\}
  \onslide<3->{&= 2^d\, 3^{\frac{dp}{p-1}}\, [w]_{A_p}^{\frac{1}{p-1}}\, \Big\| M^c_w\Big( \big(M^c_\sigma(|f|\sigma^{-1})\big)^{p-1} w^{-1} \Big) \Big\|_{L^{p/(p-1)}_w}^{\frac{1}{p-1}} \\}
  \onslide<4->{&\le 2^d\, 3^{\frac{dp}{p-1}}\, B_{\frac{p}{p-1},d}^{\frac{1}{p-1}}\, [w]_{A_p}^{\frac{1}{p-1}}\, \Big\| \big(M^c_\sigma(|f|\sigma^{-1})\big)^{p-1} w^{-1} \Big\|_{L^{p/(p-1)}_w}^{\frac{1}{p-1}} \\}
  \onslide<5->{&\le 2^d\, 3^{\frac{dp}{p-1}}\, B_{\frac{p}{p-1},d}^{\frac{1}{p-1}}\, [w]_{A_p}^{\frac{1}{p-1}}\, B_{p,d}\, \|f\|_{L^p_w}.}
\end{align*}
\onslide<6->{Choosing $D_{p,d} := 2^d\, 3^{\frac{dp}{p-1}}\, B_{\frac{p}{p-1},d}^{\frac{1}{p-1}}\, B_{p,d}$ proves (i). \hfill $\square$}
\end{frame}

%%%%%%%%%%%%%%%
\begin{frame}[t]
  \frametitle{Proof of Theorem 2 (ii) (Setup)}
\only<1>{%
\begin{block}{Outline}
Fix $\alpha < \frac{1}{p-1}$ and assume a constant $D$ such that $\|Mf\|_{L^p_w} \le D\, [w]_{A_p}^{\alpha}\, \|f\|_{L^p_w}$ for all $w \in A_p$ and all $f \in L^p_w(\mathbb{R}^d)$. Construct a family of weights and functions and derive a contradiction.
\end{block}}
\only<2-3>{%
\begin{block}{Choose}
For $0 < \delta < 1$ let
\vspace{-0.6em}
\begin{align*}
  w_\delta(x) &:= |x|^{(d-\delta)(p-1)}, \\
  f_\delta(x) &:= \sigma_\delta(x)\,\chi_{Q(0,2)}(x) = w_\delta(x)^{-\frac{1}{p-1}}\chi_{Q(0,2)}(x) = |x|^{\delta-d}\chi_{Q(0,2)}(x).
\end{align*}
\end{block}}
\only<3>{%
\begin{block}{Key identity}
For $0 < r \le 1$,
\vspace{-0.5em}
\[
  \int_{B(0,r)} f_\delta\,dy = \int_{B(0,r)} |y|^{\delta-d}\,dy = |S^{d-1}| \int_0^r s^{\delta-d}\, s^{d-1}\,ds = |S^{d-1}|\, \frac{r^\delta}{\delta} .
\]
\end{block}}
\end{frame}

%%%%%%%%%%%%%%%
\begin{frame}[t]
  \frametitle{Proof of Theorem 2 (ii) (Upper bound for $[w_\delta]_{A_p}$)}
\begin{align*}
  Q \text{ far from } 0: &\quad \bigg( \frac{1}{|Q|}\int_Q w_\delta\,dy \bigg) \bigg( \frac{1}{|Q|}\int_Q \sigma_\delta\,dy \bigg)^{p-1} \lesssim 1 \\
  \onslide<2->{Q \text{ near } 0:} &\quad \onslide<2->{\bigg( \frac{1}{|Q|}\int_Q w_\delta\,dy \bigg) \bigg( \frac{1}{|Q|}\int_Q \sigma_\delta\,dy \bigg)^{p-1} \lesssim \delta^{1-p}}
\end{align*}
\pause\pause
\begin{block}{Upper bound}
\vspace{-1.1em}
\[
  [w_\delta]_{A_p} \le K_1\, \delta^{1-p}
\]
\end{block}
\end{frame}

%%%%%%%%%%%%%%%
\begin{frame}[t]
  \frametitle{Proof of Theorem 2 (ii) (Lower bound for $Mf_\delta$)}
For $x \in Q(0,2)$ and $r := |x|_\infty$,
\begin{align*}
  Mf_\delta(x) &\ge \frac{1}{|Q(0,2r)|}\int_{B(0,r)} |y|^{\delta-d}\,dy = \frac{1}{(2r)^d}\cdot\frac{|S^{d-1}|\,r^\delta}{\delta} = \frac{|S^{d-1}|}{2^d\,\delta}\,r^{\delta-d} \\
  &\onslide<2->{\ge \frac{|S^{d-1}|}{2^d\,\delta}\,|x|^{\delta-d} = K_2\,\delta^{-1} f_\delta(x).}
\end{align*}
\pause\pause
\begin{block}{Lower bound}
\vspace{-1.1em}
\[
  \|Mf_\delta\|_{L^p_{w_\delta}} \ge K_2\, \delta^{-1}\, \|f_\delta\|_{L^p_{w_\delta}}
\]
\end{block}
\end{frame}

%%%%%%%%%%%%%%%
\begin{frame}[t]
  \frametitle{Proof of Theorem 2 (ii) (Contradiction)}
{%
\[
  \|Mf_\delta\|_{L^p_{w_\delta}} \le D\, [w_\delta]_{A_p}^{\alpha}\, \|f_\delta\|_{L^p_{w_\delta}}
\]
\uncover<2->{%
\[
  K_2\, \delta^{-1} \le \frac{\|Mf_\delta\|_{L^p_{w_\delta}}}{\|f_\delta\|_{L^p_{w_\delta}}} \le D\, [w_\delta]_{A_p}^{\alpha} \le D\, K_1^{\alpha}\, \delta^{-\alpha(p-1)}
\]}
\uncover<3->{%
\[
  \delta^{\alpha(p-1) - 1} \le \frac{D\, K_1^{\alpha}}{K_2}
\]}
\uncover<4->{%
\[
  \text{but } \delta^{\alpha(p-1) - 1} \to \infty \text{ as } \delta \to 0, \text{ contradiction}. \qquad \square
\]}}
\end{frame}

%%%%%%%%%%%% End frame
\begin{frame}[plain]
  \endpage
\end{frame}

%%%%%%%%%%%
\end{document}